\documentclass[11pt]{article}
\usepackage[T1]{fontenc}
\usepackage[margin=1.02in]{geometry}
\usepackage{amsmath,amssymb,amsthm,mathtools}
\usepackage{graphicx}
\usepackage{xcolor}
\usepackage[hidelinks]{hyperref}
\usepackage{microtype}
\newtheorem{theorem}{Theorem}
\newtheorem{proposition}{Proposition}
\newtheorem{lemma}{Lemma}
\newcommand{\R}{\mathbb R}
\newcommand{\one}{\mathbf 1}
\newcommand{\norm}[1]{\left\lVert #1\right\rVert}
\newcommand{\Afour}{\ensuremath{\mathrm{(A4)}}}
\definecolor{statusred}{RGB}{156,36,31}
\title{Stable Distance Kernels Fail the\\RKBS Admissibility Condition \Afour}
\author{fa\_banach\_001 / agent\_lane\_12}
\date{22 July 2026}
\begin{document}
\maketitle
\noindent
\fcolorbox{statusred}{statusred!6}{%
  \begin{minipage}{0.92\textwidth}
  \textbf{Status: candidate counterexample, likely valid.}
  For every $0<\gamma<1$, the kernels
  $\exp(-\norm{x-y}_p^\gamma)$ proposed in the source fail \Afour:
  the Euclidean kernel fails already in $\R^2$, and the $\ell^1$ kernel
  fails in an explicit finite dimension depending on $\gamma$.
  \end{minipage}}

\section{Source question and exact scope}

Song, Zhang, and Hickernell define an admissible kernel in
\cite[Definition~1]{SZH}.  If $x_1,\ldots,x_n,t$ are pairwise distinct,
$G=[K(x_i,x_j)]_{i,j=1}^n$, and
$k_t=(K(t,x_j))_{j=1}^n$, condition \Afour is
\begin{equation}
  \norm{G^{-1}k_t}_{\ell^1}\leq 1.
  \label{eq:A4}
\end{equation}
At the end of Section~6, on PDF page~25, the source says:
\begin{quote}
Finally, we remark from numerical experiments that the following kernels
\[
 \exp\!\left(-\norm{x-y}_{\ell^p(\mathbb N_d)}^\gamma\right),
 \qquad x,y\in\R^d,\quad \gamma\in(0,1),\quad p=1,2,
\]
seem to satisfy \Afour\ for small enough $\gamma$ and moderate $n$.
We shall leave the search of more kernels satisfying \Afour\ and its
relaxation as an open question for future study.
\end{quote}

\begin{figure}[ht]
  \centering
  \includegraphics[width=\textwidth]{figures/open_problem_crop.png}
  \caption{The stable-distance candidates and open question on page~25
  of the source paper.}
  \label{fig:source}
\end{figure}

The broad request to classify more admissible or relaxed kernels is open-ended.
The theorem below gives a full negative answer for the two displayed
multidimensional candidate families.  It does not decide their
one-dimensional versions when $0<\gamma<1$.

\section{Proof intuition}

The obstacle is visible on symmetric nodes surrounding the interpolation
target.  Every node then has the same covariance with the target, and every
row of the Gram matrix has the same sum.  Thus the interpolation coefficient
vector is constant, without inverting a large matrix.  Its $\ell^1$ norm is
the quotient of two scalar exponential sums.

At zero scale, all kernel entries tend to one and this quotient tends to one.
The first derivative records only the sum of powered distances among the
nodes.  Cross-polytopes make that derivative positive in sufficiently high
dimension for both $p=1$ and $p=2$.  In the Euclidean plane, sufficiently
large regular polygons do the same for every $\gamma>0$: their mean powered
chord length is greater than the radius power.  Hence the quotient crosses
above one at arbitrarily small positive scales and violates \eqref{eq:A4}.

\section{Counterexample theorem}

\begin{theorem}
\label{thm:main}
Fix $0<\gamma<1$.
\begin{enumerate}
\item For $p=1$, choose any integer $d$ such that
\begin{equation}
  (2d-1)2^\gamma>2d.
  \label{eq:d1}
\end{equation}
Then $K(x,y)=\exp(-\norm{x-y}_1^\gamma)$ fails \Afour\ on $\R^d$.
Consequently it fails on every $\R^{d'}$ with $d'\geq d$.
\item For $p=2$, the kernel
$K(x,y)=\exp(-\norm{x-y}_2^\gamma)$ fails \Afour\ on $\R^2$, and hence
on every $\R^d$ with $d\geq2$.
\end{enumerate}
In both cases there are violating node sets contained in every prescribed
neighborhood of the target.
\end{theorem}

The condition \eqref{eq:d1} is always achievable, since $2^\gamma>1$.
For example, the least cross-polytope dimensions selected by
\eqref{eq:d1} are $15,8,4,2$ for
$\gamma=0.05,0.10,0.25,0.50$, respectively.

\subsection{A symmetry lemma}

\begin{lemma}
\label{lem:symmetry}
Let $G$ be an invertible $n$-by-$n$ matrix with
$G\one=S\one$, where $S>0$.  If $k_t=a\one$ with $a>0$, then
\[
  G^{-1}k_t=\frac{a}{S}\one,
  \qquad
  \norm{G^{-1}k_t}_{\ell^1}=\frac{na}{S}.
\]
\end{lemma}
\begin{proof}
The displayed vector solves $Gc=k_t$.  Invertibility gives uniqueness, and
all its entries are positive, so its $\ell^1$ norm is their sum.
\end{proof}

\subsection{Cross-polytopes: both values of p}

\begin{proposition}
\label{prop:cross}
For each $p\in\{1,2\}$ and $0<\gamma<1$, the kernel
$\exp(-\norm{x-y}_p^\gamma)$ fails \Afour\ in some finite dimension.
The failure occurs on arbitrarily small cross-polytopes centered at the
target.
\end{proposition}

\begin{proof}
Put $n=2d$, take the target $t=0$, and take the nodes
\[
  x_1,\ldots,x_n\in\{\pm r e_1,\ldots,\pm r e_d\}.
\]
Write $u=r^\gamma$.  In either norm every entry of $k_0$ is
$a=e^{-u}$, and symmetry makes all Gram row sums equal.

For $p=1$, every two distinct nodes have distance $2r$.  With
$b=e^{-2^\gamma u}$,
\[
  G=(1-b)I+b\one\one^{T},\qquad
  S_1(u)=1+(n-1)e^{-2^\gamma u}.
\]
The eigenvalues are $1-b>0$ and $S_1(u)>0$, so $G$ is invertible.
By Lemma~\ref{lem:symmetry}, \Afour\ fails exactly when
\[
 F_1(u):=ne^{-u}-1-(n-1)e^{-2^\gamma u}>0.
\]
Now $F_1(0)=0$ and
\[
 F_1'(0)=(n-1)2^\gamma-n.
\]
Choose $n=2d$ satisfying \eqref{eq:d1}.  Then $F_1'(0)>0$, and therefore
$F_1(u)>0$ for all sufficiently small $u>0$.

For $p=2$, each node has one antipode at distance $2r$ and $n-2$ other
nodes at distance $\sqrt2\,r$.  Hence
\[
 S_2(u)=1+e^{-2^\gamma u}
          +(n-2)e^{-2^{\gamma/2}u}.
\]
For completeness, this Gram matrix is invertible without appealing to a
general kernel theorem.  If $P$ exchanges each node with its antipode and
$b=e^{-2^\gamma u}$, $c=e^{-2^{\gamma/2}u}$, then
\[
 G=(1-c)I+(b-c)P+c\one\one^T.
\]
Its eigenvalues are $S_2(u)$, $1-b$, and $1+b-2c$.  The last is positive:
\[
 \frac{1+b}{2}>\sqrt b
 =e^{-2^{\gamma-1}u}>e^{-2^{\gamma/2}u}=c,
\]
where $\gamma<2$ gives $2^{\gamma-1}<2^{\gamma/2}$.

Again Lemma~\ref{lem:symmetry} reduces failure to $F_2(u)>0$, where
\[
 F_2(u):=ne^{-u}-S_2(u),\qquad
 F_2'(0)=2^\gamma+(n-2)2^{\gamma/2}-n.
\]
Since $2^{\gamma/2}>1$, the derivative is positive for all sufficiently
large even $n$.  Then $F_2(u)>0$ for every sufficiently small $u>0$.
As $r=u^{1/\gamma}$ can be arbitrarily small, so can the configurations.
\end{proof}

\subsection{Regular polygons: Euclidean failure in dimension two}

\begin{proposition}
\label{prop:circle}
For every $\gamma>0$ for which
$K_\gamma(x,y)=e^{-\norm{x-y}_2^\gamma}$ is strictly positive definite
(in particular, every $0<\gamma<1$), the kernel fails \Afour\ on $\R^2$.
\end{proposition}

\begin{proof}
For an integer $n\geq3$, let the target be $0$ and put
\[
 x_j=r\bigl(\cos(2\pi j/n),\sin(2\pi j/n)\bigr),
 \qquad j=0,\ldots,n-1.
\]
The nodes are distinct and the stable radial kernel is strictly positive
definite for $0<\gamma<2$ by Schoenberg's theorem
\cite{Schoenberg}; hence their Gram matrix is invertible.

Set $u=r^\gamma$ and
\[
 \delta_j=2\sin(\pi j/n),\qquad
 S_n(u)=\sum_{j=0}^{n-1}e^{-u\delta_j^\gamma}.
\]
The Gram matrix is circulant with row sum $S_n(u)$, while
$k_0=e^{-u}\one$.  Lemma~\ref{lem:symmetry} gives
\begin{equation}
 \norm{G^{-1}k_0}_{\ell^1}
   =\frac{ne^{-u}}{S_n(u)}.
 \label{eq:circlequot}
\end{equation}

We first choose $n$.  The continuous Riemann sums satisfy
\begin{equation}
 \frac1n\sum_{j=0}^{n-1}\delta_j^\gamma
 \longrightarrow
 I_\gamma:=\int_0^1(2\sin(\pi s))^\gamma\,ds.
 \label{eq:riemann}
\end{equation}
Moreover $I_\gamma>1$.  Indeed,
\[
 \int_0^1\log(2\sin(\pi s))\,ds=0,
\]
because $\int_0^\pi\log(\sin t)\,dt=-\pi\log2$.  The logarithm is
integrable and nonconstant, so strict Jensen gives
\[
 I_\gamma
 =\int_0^1\exp\!\bigl(\gamma\log(2\sin(\pi s))\bigr)\,ds
 >
 \exp\!\left(\gamma\int_0^1\log(2\sin(\pi s))\,ds\right)=1.
\]
By \eqref{eq:riemann}, some $n$ therefore satisfies
\begin{equation}
  \sum_{j=0}^{n-1}\delta_j^\gamma>n.
  \label{eq:chord}
\end{equation}

For this $n$, define $F_n(u)=ne^{-u}-S_n(u)$.  Then
\[
 F_n(0)=0,\qquad
 F_n'(0)=-n+\sum_{j=0}^{n-1}\delta_j^\gamma>0
\]
by \eqref{eq:chord}.  Thus $F_n(u)>0$ for all sufficiently small
$u>0$.  Equation~\eqref{eq:circlequot} is then strictly greater than one,
which violates \Afour.  Taking $u\downarrow0$ places the polygon in an
arbitrarily small disk.
\end{proof}

\begin{proof}[Proof of Theorem~\ref{thm:main}]
The $\ell^1$ assertion follows from the first half of
Proposition~\ref{prop:cross}; embedding the same finite configuration by
zero-padding gives failure in every higher dimension.  The Euclidean
assertion is Proposition~\ref{prop:circle}, and the same embedding gives all
$d\geq2$.  Both constructions allow $r$ to be arbitrarily small.
\end{proof}

\section{Verification}

The proof above is analytic.  The deterministic script
\path{code/verify_symmetric_counterexamples.py} independently:
\begin{itemize}
\item finds cross-polytope dimensions and checks the scalar quotients at
80-digit precision for six values of $\gamma$ from $0.05$ to $0.95$;
\item finds regular-polygon sizes and checks their powered-chord derivatives;
\item directly builds and solves a $6$-by-$6$ Euclidean Gram system for
$\gamma=1/2$, rather than using the quotient formula.
\end{itemize}
The direct solve reports minimum eigenvalue
$7.79\cdot10^{-3}$, infinity-norm residual
$1.11\cdot10^{-16}$, coefficient spread
$1.69\cdot10^{-15}$, and \Afour\ norm $1.000065821370$.
The run ends with
\texttt{ALL SYMMETRIC COUNTEREXAMPLE CHECKS PASSED}.
These checks are not used as proof.

\section{Limitations and bounded novelty check}

The theorem settles the two named candidates negatively as multidimensional
\Afour\ kernels.  It does not classify all admissible kernels, decide either
kernel on $\R$, or give optimal growth for the relaxed Lebesgue constants.
The source's finite numerical observation for moderate $n$ is not itself
false: the theorem shows why such experiments cannot establish the universal
quantifiers in \Afour.

On 22 July 2026, the run's registry, solution, attempt, and proof-gap indexes
were searched by arXiv id, exact title, the phrases admissible kernel,
condition (A4), stable distance kernel, and the displayed exponential formula.
Bounded web/arXiv searches used the exact source sentence, title/citation
combinations, and the formula together with Lebesgue constant and A4.
They found the source and later papers constructing other admissible RKBS
kernels, including arXiv:1901.01036 and arXiv:1903.00819, but no prior
resolution of this named stable-distance family.  This is a bounded search,
not an exhaustive novelty claim.

\section{Human-review recommendation}

This packet is suitable for expert review as a candidate counterexample.
The highest-value checks are the derivative signs in
Proposition~\ref{prop:cross}, strict Jensen and the Riemann-sum quantifier in
Proposition~\ref{prop:circle}, and the use of strict positive definiteness for
the planar Gram matrix.  Until review, retain the status
\emph{candidate counterexample, likely valid}.

\begin{thebibliography}{9}
\bibitem{SZH}
G. Song, H. Zhang, and F. J. Hickernell,
\emph{Reproducing Kernel Banach Spaces with the $\ell^1$ Norm},
Appl. Comput. Harmon. Anal. \textbf{34} (2013), 96--116;
arXiv:1101.4388.

\bibitem{Schoenberg}
I. J. Schoenberg,
\emph{Metric spaces and positive definite functions},
Trans. Amer. Math. Soc. \textbf{44} (1938), 522--536.
\end{thebibliography}
\end{document}
